\documentclass{jsarticle}
\usepackage[dvipdfmx,final]{graphicx}
\usepackage{amsmath}
\usepackage{amssymb}
\begin{document}
\title{M theory and time measure potential equation \\
and qunatum level of space ideality thoerem}
\author{Masaaki Yamaguchi}
\date{}
\maketitle
\newcommand{\timebow}{%
	\begin{picture}(5,5)
		\put(0,0){$\text{\rotatebox{90}{$\bowtie$}}$}
	\end{picture} %
	}
   \newcommand{\variint}{%
	\begin{picture}(15,30)
		\put(0,0){
			$ \iint $}
		\put(0,5){\line(15,0){15}}
	\end{picture} %
}

\newcommand{\variiint}{%
	\begin{picture}(15,15)
		\put(0,0){
			$ \iiint $}
		\put(0,5){\line(15,0){20}}
	\end{picture} %
}



 \newcommand{\alearletter}{%
	 \begin{picture}(10,20)
		 \put(0,0){
			 $ \square $}
		 \put(0,0){\line(1,1){10}}
	 \end{picture} %
	 }


       時間計量による、指数定理における擬微分と擬積分多様体を微小変量によって、
   大域的トポロジーが成されられる。ガンマ関数における大域的多様体と、
   ヒッグス方程式における大域的多様体の計算式を以下に、載せておく。
   $$ {T^{\mu\nu}}^{T^{\mu\nu}} = \int T^{\mu\nu} \rotatebox{90}{$\bowtie$}_m $$
   $$   {T^{\mu\nu}}^{{T^{\mu\nu}}^{'}}=\int T^{\mu\nu}dx_m  $$
      $$ {\Gamma^{\gamma}}^{'} = {d \over {d \gamma}}\Gamma $$
   $$ \Gamma^{\gamma} = \int \Gamma dx_m $$
   $$ \rotatebox{90}{$\bowtie$} = \int x^x dx_m $$


     ヒッグス場方程式が、相加相乗平均方程式であり、大域的多様体でもある。
  $$ {d \over {df}}F(x,y)=m(x) $$
  $$ = {d \over {d m(x)}}M = \partial (\sigma(M)) $$
  $$=  {d \over {d \timebow}}Z^{'}(\zeta)dx_m $$
  $$ ||ds^2|| = \nabla_{i}\nabla_{j}\int \nabla M(\sigma)d\eta $$
  M理論の指数定理による擬微分と擬積分多様体が、ヘルマンダー作用素から
  のノルム空間となり、
  $$ = M^{{m(x)^{'}}} $$
  M理論のニュートン方式の大域的微分多様体になる。
  $$ ||ds^2|| = \int_M [d(\partial) + d(\sigma)]dI_m $$
  それが、D-braneのM理論へと導かれる。
  $$ {d \over {df}}(e^{-f} + e^{f}) = {{d \over {d R^{'}}}}R = d[I_m] $$
  結局は、Jones多項式における、大域的多様体としての、虚数微分方程式となる。
  リッチ・フロー方程式による、大域的多様体の擬微分と擬積分多様体の
  計算式が、ニュートン方式による、擬積分と擬微分多様体となり、
  時間計量が、相対性による、測量がなくなり、空間計量となり、
  空間概念の量子化となる。

\end{document}
